\begin{definition}\cite{sim:vol1,sim:vol3}
	A finite dimensional algebra $\Lambda$ is said to be of \textbf{finite representation type} (or \textbf{representation-finite}) if the category $\Modf{\Lambda}$ has a finite number of isomorphism classes of indecomposable modules. Otherwise it is said to be of \textbf{infinite representation type} (or \textbf{representation-infinite}).\smallskip
	
	Moreover $\Lambda$ is said to be:
	\begin{itemize}
		\item of \textbf{tame representation type} (shortly \textbf{tame}) if there is a classification of the isomorphism classes of the indecomposable modules in $\Modf{\Lambda}$ in the sense
		that, for each integer $d \geq 1$, the indecomposable modules in $\Modf{\Lambda}$ of dimension $d$ form at most finitely many one-parameter families.
		\item of \textbf{wild representation type} (shortly \textbf{wild}) if, for any finite dimensional $\K$-algebra $A$, there exists an embedding functor $T\colon \Modf{A}\to \Modf{\Lambda}$.
	\end{itemize}
\end{definition}

%%%%%%%%%%%%%%%%%%%%%%%%%%%%%%%%%%%%%%%%%%%%%%%%%%%%%%%%%%%%%%%%%%%%

\subsection{Tame concealed algebras}
We will work with tame concealed algebras. We recall their definition and some useful properties.

\begin{definition}[{c.f.~\cite[Ch.VIII 3.1, 4.6]{sim:vol1}}, {\cite[Ch.XI.3]{sim:vol2}}]
	Let $Q$ be an acyclic quiver. A finite-dimensional $\K$-algebra $\Lambda$ is said to be \textbf{tilted of type $Q$} if it is isomorphic to $\End_{\K Q}(T)$ for some tilting $\K Q$-module $T$. If $T$ is also postprojective, $\Lambda$ is said to be \textbf{concealed of type $Q$}.
	
	We say that a concealed $\K$-algebra of type $Q$ is \textbf{tame concealed} (equiv.~\textbf{concealed of Euclidean Type}) if $Q$ is a tame and representation-infinite quiver (i.e.~an Euclidean quiver).
\end{definition}

We have the following result, justifying the name ``tame concealed'':
\begin{theorem}\label{thm:tameconcealed}
	A tame concealed algebra is representation-infinite, tame and $\tau$-tilting infinite.
\end{theorem}
\begin{proof}
	It follows directly from the classification of tame concealed algebras that they are representation-infinite (see \cite[Theorem 1]{hap:min}, \cite[Section 4.3, Proposition 7]{ringel2006tame}). They are also tame by \cite[Ch.XIX, Theorem 3.14]{sim:vol3}, and $\tau$-tilting infinite (see \cite[Lemma 2.6]{borve:tau}).
\end{proof}


%%%%%%%%%%%%%%%%%%%%%%%%%%%%%%%%%%%%%%%%%%%%%%%%%%%%%%%%%%%%%%%%%%%

\subsection{Tits form}
We will need the notion of the Tits form of a bound quiver. We present here some definitions and results that will help us in the following.

\begin{definition}[{c.f.~\cite[Section 1.0]{ringel2006tame}}]
	An \textbf{integral quadratic form} is a polynomial $q=q(x_1,\dots,x_n)$ in $n$ variables with integral coefficients $q_{ij}\in \Z$ of the form
	\[q=q(x_1,\dots,x_n)=\sum_{i=1}^{n} x_i^2 + \sum_{i<j} q_{ij} x_i x_j\]
	We obtain a function $q\colon\Z^n \to \Z$ evaluating the polynomial $q$ at $n$-tuples of integral numbers. 
	
	We can endow $\Z^n$ with a partial order defined component-wise. Given $z = (z_1,\dots,z_n) \in \Z^n$, it is said to be:
	\begin{itemize}
		\item \textbf{positive}, denoted by $z > 0$, if $z\neq 0$ and $z_i\geq0$ for all $i$;
		\item \textbf{sincere} if $z_i\neq 0$ for all $i$.
	\end{itemize}
	
	An integral quadratic form $q$ in $n$ variables is said to be \textbf{positive semi-definite} if $q(z)\geq 0$ for all $z\in\Z^n$.
	
	For a positive semi-definite quadratic form $q$, we can define the \textbf{radical of $q$} as 
	\[\Rad(q)\defeq \{z\in \Z^n\sth q(z) = 0\}\]
	It is a subgroup of $\Z^n$ and its rank is called the \textbf{radical rank of $q$}.
	The elements $z\in\Rad(q)$ are called \textbf{radical vectors}.
	
	In case $0$ is the only radical	vector of $q$ or, equivalently, the radical rank of $q$ is $0$, the form $q$ is said to be \textbf{positive definite}.
\end{definition}

\begin{definition}[{c.f.~\cite[Section 1.1]{ringel2006tame}}]\label{def:quadmatrix}
	Given $q=q(x_1,\dots,x_n)=\sum_{i=1}^{n} x_i^2 + \sum_{i<j} q_{ij} x_i x_j$ an integral quadratic form in $n$ variables, we define the corresponding symmetric bilinear form $\beta\colon \Z^n \times \Z^n \to \Z$ as $\beta(x,y) \defeq x^t B y$ with $B$ the symmetric $n\times n$-matrix $B=\frac{1}{2}(q_{ij})_{1 \leq i,j \leq n}$, with $q_{ij}=q_{ji}$ for $i > j$, and $q_{ii} = 2$ for all $i$.	
\end{definition} 

\begin{lemma}\label{lem:radpossemi}
	Let $q$ be an integral quadratic form, $\beta$ the corresponding symmetric bilinear form with matrix $B$. If $q$ is positive semi-definite then we have that
	\[\Rad(q)=\ker(B)\]
\end{lemma}
\begin{proof}
	($\supseteq$) If $z\in\Z^n$ is in $\ker(B)$ (i.e.~$Bz=0$) then $q(z)=\beta(z,z)=z^t B z = 0$. Thus $z$ is in the radical $\Rad(q)$ of $q$.
	
	($\subseteq$) Let $z$ be in the radical $\Rad(q)$ of $q$, so $q(z)=0$. Since $q$ is positive semi-definite, we have that $\beta$ is a positive semi-definite bilinear form, i.e.~$B$ is a positive semi-definite matrix (all its eigenvalues are $\geq0$). By Cholesky factorisation (see \cite[Corollary 7.2.9]{horn2012matrix}) we can write $B= L^t L$ with $L$ a real lower triangular matrix with non-negative diagonal entries. Therefore
	\[
	0=q(z)=\beta(z,z)=z^t B z= z^t L^t Lz=(Lz)^t Lz = \lVert Lz\rVert ^2
	\]
	with $\lVert\mhyphen\rVert$ being the (real) Euclidean norm ($\lVert x\rVert^2 =x^t x= \sum_{i=1}^{n} x_i^2$). Then $\lVert Lz\rVert ^2=0$, which implies that $Lz=0$. Thus $Bz=L^tLz=0$, so $z\in\ker(B)$.
\end{proof}

\begin{definition}[{c.f.~\cite[Section 1.0]{ringel2006tame}}]
	A quadratic form $q$ in $n$ variables is said to be \textbf{weakly positive} if $q(z) > 0$ for all positive $z\in\Z^n$.
	
	For any $1\leq t\leq n$, we denote by $q^t= q(z_1,\dots, z_{t-1},0,z_{t+1},\dots,z_n)$ the restriction of $q$ to the hyper-plane defined by $z_t=0$.
	
	An integral quadratic form $q$ in $n\geq 3$ variables is said to be \textbf{critical} if $q$ is not weakly positive but all the
	forms $q^t$, $1\leq t\leq n$, are weakly positive.
\end{definition}

We have the following result due to Ovsienko:
\begin{theorem}[{c.f.~\cite[Section 1.0, Theorem 2]{ringel2006tame}}]\label{thm:criticalform}
	A critical quadratic form is positive semi-definite	with radical rank $1$, and with a sincere positive radical vector.
\end{theorem}

Now we define some integral quadratic forms over algebras and quivers.

\begin{definition}[{c.f.~\cite[Section 2.2]{bongartz83quad}}]
	Let $\Lambda$ be a finite dimensional $\K$-algebra with finite global dimension ($\gldim{\Lambda}\leq m$), $\Lambda\cong \palgebra{}/I$ for some finite quiver $Q$ (with $n= \card{Q_0}$) and admissible ideal $I$. We define the \textbf{Euler characteristic} of $\Lambda$ as the bilinear form 
	\[
	\langle \dimvec{M},\dimvec{M'}\rangle \defeq \sum_{i\geq0} (-1)^i \dim_\K \Ext_{\Lambda}^i(M,M')  
	\]
	where $M$, $M'$ are finitely generated $\Lambda$-modules. It is well defined since $\gldim{\Lambda}\leq m$, so $\Ext_{\Lambda}^i(M,M')=0$ for $i\geq m+1$. We define the \textbf{Euler quadratic form} of $\Lambda$ as
	\[\chi_\Lambda(\dimvec{M})\defeq \langle \dimvec{M},\dimvec{M}\rangle\]
\end{definition}

\begin{definition}[{c.f.~\cite[Definition 2.1]{bongartz83quad}}]
	Let $(Q,I)$ be an arbitrary bound quiver with $n=\card{Q_0}$ and $Q$ acyclic. The \textbf{Tits quadratic form} of $\Lambda=\palgebra{}/I$ is the integral quadratic form $q_\Lambda \colon \Z^{Q_0} \to \Z$ defined by the formula:
	\[
	\begin{aligned}
		q_\Lambda(x)\defeq& \sum_{i\in Q_0} x_i^2 - \sum_{\alpha\in Q_1} x_{s(\alpha)} x_{t(\alpha)} + \sum_{i,j \in Q_0} r_{ij} x_i x_j \\
		=& \sum_{i\in Q_0} x_i^2 - \sum_{i,j\in Q_0} d_{ij} x_i x_j + \sum_{i,j \in Q_0} r_{ij} x_i x_j
	\end{aligned}
	\]
	where $d_{ij}$ is the number of arrows from $i$ to $j$ and $r_{ij} = \card{R\cap \varepsilon_j \palgebra{} \varepsilon_i}$, with $R$ a minimal set of relations which generates $I$ and $\varepsilon_j \palgebra{} \varepsilon_i$ the vector space spanned by all the paths starting at $i$ and ending in $j$ ($\varepsilon_\ell$ denotes the lazy path at vertex $\ell$) .
\end{definition}

These two integral quadratic forms may be different in general, but they coincide with the following hypothesis:

\begin{theorem}[{c.f.~\cite[Proposition 2.2]{bongartz83quad}}]\label{thm:eulertits}
	Let $\Lambda =\palgebra{}/I$ be a basic finite-dimensional $\K$-algebra  with $\gldim{\Lambda}\leq 2$ and $Q$ acyclic.  Then $\chi_\Lambda$ and $q_\Lambda$ coincide.
\end{theorem}

%%%%%%%%%%%%%%%%%%%%%%%%%%%%%%%%%%%%%%%%%%%%%%%%%%%%

\subsection{Incidence algebras}

We will now present some notions related to the representation theory of a finite poset $P$.

\begin{definition}[{c.f.~\cite[Definition 1.1]{borve:tau}}]\label{def:incidencealgebra}
	Let $P$ be a finite poset. Its \textbf{Hasse quiver} $\HasseQuiver{P}$ is defined as the quiver with the elements in $P$ as vertices and arrows $x\rightarrow y$ if $x\leq y$ in $P$ and no element lies strictly between $x$ and $y$.
	
	We define the \textbf{incidence algebra} $\incidencealgebra{P}$ of a finite poset $P$ as the bound path algebra $\palgebra{P} / I_P$, where the ideal $I_P$ of the path $\K$-algebra $\palgebra{P}$ is generated by the relations $\{p_1-p_2\}_{p_1,p_2}$, indexing over all pairs of paths with the same start and endpoints.
\end{definition}

 In the following, (finite-dimensional) left $\incidencealgebra{P}$-modules will also be considered as (finite-dimensional) representations of the bound quiver $(\HasseQuiver{P}, I_P)$, namely $\Mod{\incidencealgebra{P}}\cong \Rep{\HasseQuiver{P}, I_P}$ (resp.~$\Modf{\incidencealgebra{P}}\cong \Repf{\HasseQuiver{P}, I_P}$).

\begin{remark}
	We would like to bring to the reader's attention that we will be working with the classical definition of representations of the bound quiver $(\HasseQuiver{P}, I_P)$.
	
	Other authors used the term \emph{representation} of a poset $P$ (or also \emph{$P$-space}) to denote a $\K$-vector space $V$ with an order-preserving map from $P$ into the poset of subspaces of $V$ (see for example \cite{nazarovaroiter:representations, gabriel:representations}). The maps in such representations of the bound quiver $(\HasseQuiver{P},I_P)$ are all injective. This property is not required in the following.
\end{remark}

\begin{definition}
	A finite poset $P$ is said to be \textbf{representation-infinite} if the incidence algebra $\incidencealgebra{P}$ is representation-infinite. The definition of a representation-finite, tame or wild poset is analogous.
\end{definition}


\subsubsection{Critical posets and reduction techniques}
 
Loupias \cite{lou:ind,lou:repr} devised two reduction techniques of posets that proved to be very useful to study the representation theory of incidence algebras.
\begin{definition}[{\cite{lou:ind}, c.f.~\cite[Section 2.3]{borve:tau}}]\label{def:reduction}
	Let $P$ be a finite poset.
	\begin{itemize}
		\itemsep0em
		
		\item A subposet $P'$ of $P$ (with the induced order and the canonical order-preserving map $s\colon P'\to P$) is also called a \textbf{subtraction} of $P$.
		
		\item A poset $P'$ is a \textbf{contraction} (or a \textbf{contracted poset}) of $P$ if there exists a surjective order-preserving morphism $c\colon P\rightarrow P'$ (also called \textbf{contraction}) with connected fibers (i.e.~$c^{-1}(x)$ is connected for every $x\in P'$).
		
		\noindent If there exists exactly one element $x\in P'$ such that $\vert c^{-1}(x)\vert=2$ and $|c^{-1}(y)|=1$ for all $y\in P'\setminus \{x\}$, the contraction $c\colon P\to P'$ is called \textbf{elementary contraction}.
		
		\item A poset $P$ \textbf{can be reduced to $P'$} if $P'$ can be obtained from $P$ using a finite number of subtractions and/or contractions.
	\end{itemize}
\end{definition}
Informally, an elementary contraction identifies two adjacent elements in $P$. Therefore, any arrow in the Hasse quiver $P$ corresponds to an elementary contraction. It is possible to check that every contraction can be written as a composite of elementary contractions (c.f.~\cite[Lemma 2.29]{tacchetti2025masterthesis}).

If a poset $P$ can be reduced to $P'$, we may assume that each step of the reduction process either removes a single vertex or contracts an arrow in the Hasse quiver.

\begin{definition}[{\cite[Section 1]{lou:ind}}]\label{def:criticalposet}
	A poset $P$ is said to be \textbf{critical} (or $\incidencealgebra{P}$ is said to be \textbf{minimally representation-infinite}) if it is representation-infinite and all of its	proper subposets and non-trivial contractions have representation-finite incidence $\K$-algebra.
\end{definition}

\begin{lemma}\label{lemma:infinitereduction}
	Any representation-infinite poset can be reduced to a critical poset.
\end{lemma}
\begin{proof}
	Let $P$ be a representation-infinite poset. If all of its proper subposets and non-trivial contractions have representation-finite incidence $\K$-algebras, it is critical by Definition \ref{def:criticalposet}.
	Otherwise, $P$ will have at least one proper subposet or non-trivial contraction $P'$ that is representation-infinite. The cardinality of $P'$ is strictly less than the cardinality of $P$. Therefore by induction on the cardinality of the poset, $P'$ can be reduced to a critical poset, thus the same holds for $P$, since we can consider the reduction techniques used from $P$ to $P'$ and from $P'$ to a critical poset in order to reduce $P$ to a critical poset.
\end{proof}

\captionsetup[table]{skip=10pt}
\begin{table}[ht!]
	\centering
	\begin{tabular}{cc|cc|cc}
		$\qAtilde_3$& \posetAthreetilde{} & $\qDtilde_4$  & \posetDfourtilde{} & $\qEtilde_6$& \posetEsixtilde{} \\[3pt] \hline
		&&&&& \\[-6pt]
		$\qEtilde_7$ & \posetEseventilde{} & $\qEtilde_8$& \posetEeighttilde{} & $\qR_1$ & \posetRone{} \\[3pt]  \hline
		&&&&& \\[-6pt]
		$\qR_2$& \posetRtwo{} & $\qR_3$ & \posetRthree{} & $\qR_4$& \posetRfour{} \\[3pt] \hline
		&&&&& \\[-10pt]
		$\qR_5$ & \posetRfive{} &$\qR_6$ & \posetRsix{} & $\qR_7$ &  \posetRseven{} \\[3pt] 
	\end{tabular}
	\caption{Critical posets are given by these bound Hasse quivers (and their opposites). Dashed lines indicate relations. Undirected edges in the quiver indicate that the arrows may go in either direction.}
	\label{tab:Loupias frames}
\end{table}

Loupias showed that the critical posets are exactly the posets in Table \ref{tab:Loupias frames} or the opposites of those (see \cite[Theorem 1.4]{lou:ind}) and that they characterise representation-infinite posets through the reduction process (see \cite[Theorem 1.5]{lou:ind}). Therefore we can summarise these results with the following:

\begin{theorem}[{see \cite[Theorem 1.4, 1.5]{lou:ind}}] \label{thm:Lou75concealed}
	A finite poset $P$ is representation-infinite if and only if it can be reduced to a critical poset $C$, i.e. to a poset whose bound Hasse quiver (or its opposite) is in Table \ref{tab:Loupias frames}. 
\end{theorem}


\begin{lemma}\label{lemma:postameconc}
	Let $C$ be a critical poset. Then $\incidencealgebra{C}$ is tame concealed.
	In particular it is representation-infinite, tame and $\tau$-tilting infinite.
\end{lemma}
\begin{proof}
	Critical posets are given by bound Hasse quivers quivers in Table \ref{tab:Loupias frames} (and their opposites). Such quivers appear on Happel-Vossieck’s list of tame concealed algebras \cite{hap:min}. The second part of the statement follows immediately from Theorem \ref{thm:tameconcealed}.
\end{proof}

Thanks to this result, the incidence algebra $\incidencealgebra{C}$ of a critical poset $C$ is in particular tilted of Euclidean type (by definition), so the following result holds:

\begin{theorem}\label{thm:titscritical}
	Let $C$ be a critical poset, $\incidencealgebra{C}$ its incidence algebra.
	\begin{enumerate}
		\item[$(1)$] $\gldim{\incidencealgebra{C}}\leq 2$ and any indecomposable $X\in \Modf{\incidencealgebra{C}}$ has projective dimension or injective dimension less or equal than one.
		
		\item[$(2)$] $\HasseQuiver{C}$ is acyclic.
		
		\item[$(3)$] The Tits form $q_{\incidencealgebra{C}}$ (also denoted by $q_C$ in the following) is critical. In particular it is positive semi-definite with radical rank $1$ and with a sincere positive radical vector.
	\end{enumerate} 
\end{theorem}
\begin{proof}
	$(1)$ Follows from \cite[Ch.VIII, Lemma 3.2]{sim:vol1} since $\incidencealgebra{C}$ is tilted of Euclidean type.
	
	$(2)$ The quiver associated to a tilted algebra is acyclic (see \cite[Ch.VIII, Corollary 3.4]{sim:vol1}).
	
	$(3)$ By the proof of \cite[Proposition 4.3.7]{ringel2006tame}, tame concealed algebras have a critical Euler quadratic form, therefore $\chi_{\incidencealgebra{C}}$ (also denoted by $\chi_C$) is critical. But now, since $(1)$ and $(2)$ hold, we can apply Theorem \ref{thm:eulertits} and we have that $\chi_C=q_C$. Therefore the Tits form $q_C$ is critical and by Theorem \ref{thm:criticalform} we conclude that it is positive semi-definite with radical rank $1$ and with a sincere positive radical vector.
\end{proof}







